% Generated by tools/edn_latex.py from reports/edn.md.
% Do not edit: change reports/edn.md and run `make edn`.
\ednentry{
  id = {EDN-68},
  title = {The Great Expectations store is a cached build product, the suites are generated from the contract, and each rule sits on the frame where it can hold},
  date = {2026-10-05},
  milestone = {M3: Quality Assurance},
  topic = {Data Validation, Testing Strategy, Reproducibility},
  participants = {@lukas2510, decided by the repository owner},
  decision = {\texttt{configure\_\allowbreak{}gx} rebuilds the Great Expectations file context under \texttt{gx/\allowbreak{}} from scratch on every run, from \texttt{recommenditos/\allowbreak{}data/\allowbreak{}gx\_\allowbreak{}context\_\allowbreak{}configuration.\allowbreak{}py}, \texttt{schema.\allowbreak{}py} and \texttt{params.\allowbreak{}yaml}, and \texttt{gx/\allowbreak{}} is that stage's DVC-cached output: gitignored, pushed, and a dependency of \texttt{validate-\allowbreak{}data}. The context holds two dataframe assets, \texttt{raw} and \texttt{interim}, with one suite, validation definition and checkpoint each. The suites' column lists, dtypes and non-null checks are generated from \texttt{schema.\allowbreak{}py} rather than written by hand. Every rule of issue \#25 sits on the frame where it can hold: \texttt{registration\_\allowbreak{}date} at or before the reference date hard on the interim frame and with \texttt{validate.\allowbreak{}raw\_\allowbreak{}mostly} on the raw one, evaluated on the parsed date; the training price range on the interim frame only; the FR-03 mileage range as a tolerant check on both frames (EDN-72); no agreement rule between \texttt{is\_\allowbreak{}used} and \texttt{offer\_\allowbreak{}type}; and no supported-make rule at all, because \texttt{split} computes the list after this stage and the downstream stages apply it (EDN-48). No suite can pass vacuously: both require \texttt{validate.\allowbreak{}min\_\allowbreak{}rows} rows and the columns their bounds read filled up to \texttt{validate.\allowbreak{}filled\_\allowbreak{}mostly}, and \texttt{validate-\allowbreak{}data} refuses an empty suite and a run that returns fewer results than its suite holds. A failed expectation fails the stage after it has written a git-tracked summary without timestamps and the Data Docs, both under \texttt{reports/\allowbreak{}data-\allowbreak{}validation/\allowbreak{}}; a frame that breaks its contract is not handed to its suite.},
  rationale = {\emph{Alternatives considered:}
\begin{itemize}[leftmargin=*, topsep=2pt]
\item \textbf{Option A: commit \texttt{gx/\allowbreak{}} as configuration, as the course demo does.}
\newline Pros: the suites are readable as JSON on GitHub and the report can link them; it is the layout the course teaches.
\newline Cons: Great Expectations writes a fresh UUID into every object it saves, so measured on 1.23.2, two builds of the same code differ in 6 of the 9 files a context holds, and a second build in place still rewrites 3 of them. Every rerun of the stage would be a diff, two branches that both rebuilt the store would conflict, and the committed JSON says less readably what the module already says. The demo's Parquet assets also record the absolute path of its author's checkout, so its store works on one machine only.
\item \textbf{Option B: commit \texttt{gx/\allowbreak{}}, and pin every id so a rebuild is byte-identical.}
\newline Pros: committed and stable at once, and a test could assert the committed store is the one the code generates.
\newline Cons: Great Expectations overwrites an id passed to a suite, a validation definition or a checkpoint when it saves one (measured), so pinning means rewriting the files after it has written them. That is a post-processor against an internal format, and it would break silently on an upgrade.
\item \textbf{Option C: no store; build an ephemeral context inside \texttt{validate-\allowbreak{}data}.}
\newline Pros: the least code and nothing to version.
\newline Cons: drops the \texttt{configure\_\allowbreak{}gx} stage that \#25 and the course demo prescribe, leaves the suites with no inspectable form at all, and the Data Docs need stored results anyway.
\item \textbf{Option D (chosen): the store is \texttt{configure\_\allowbreak{}gx}'s cached DVC output, rebuilt from scratch.}
\newline Pros: no UUID noise in Git; \texttt{dvc pull} restores the exact store a validation ran against; declaring it as an output is what puts \texttt{configure\_\allowbreak{}gx} in the DAG, which the demo's stage is not; a rebuild cannot keep a suite or an expectation the module no longer defines, while the demo's script fails outright on its second run because the data source already exists.
\newline Cons: the suites are not readable on GitHub, so a reviewer reads the module that builds them, and the Data Docs render them after \texttt{dvc pull}. A rebuild changes the output hash even when no rule changed, so \texttt{validate-\allowbreak{}data} reruns whenever \texttt{configure\_\allowbreak{}gx} does -- which is only when the module, the contract or one of its parameters changed.
\end{itemize}
\par
\emph{Why:} The store question was decided by measurement rather than taste: an artefact that changes on every rebuild with no change in meaning is a build product, and this project already treats build products as DVC outputs (\texttt{models/\allowbreak{}}, EDN-56). The other choices follow from what the tool can and cannot do, each found by running it on our frames before building on it (\texttt{reports/\allowbreak{}analysis/\allowbreak{}gx\_\allowbreak{}probe.\allowbreak{}py}):
\par
\begin{itemize}[leftmargin=*, topsep=2pt]
\item \textbf{Dataframe assets, one checkpoint per frame.} A dataframe asset records no path, so the store works in every clone, and \texttt{validate-\allowbreak{}data} validates the frame its contract has just accepted instead of reading the 215 MB file twice. \texttt{Checkpoint.\allowbreak{}run} takes one set of batch parameters for all its validation definitions, so two frames need two checkpoints. The Data Docs are rendered once per stage run rather than by a checkpoint action, which would render the whole site once per frame.
\item \textbf{Results and Data Docs outside \texttt{gx/\allowbreak{}}.} They go to \texttt{reports/\allowbreak{}data-\allowbreak{}validation/\allowbreak{}results/\allowbreak{}} and \texttt{data-\allowbreak{}docs/\allowbreak{}}, \texttt{validate-\allowbreak{}data}'s own outputs. Written inside \texttt{gx/\allowbreak{}}, every validation would modify \texttt{configure\_\allowbreak{}gx}'s output and \texttt{dvc status} would report that stage changed after every run.
\item \textbf{Generated contract expectations.} \#25 puts the column schema and types in the suite; generating them from \texttt{schema.\allowbreak{}py} keeps it the one place the contract is written, and makes the Data Docs a complete statement of each frame. They pass whenever they are reached, because \texttt{Schema.\allowbreak{}validate} runs first, and they are weaker than it: Great Expectations compares only a column's scalar type, so it cannot tell nullable, extension or \texttt{object} dtypes from their counterparts. Measured, \texttt{datetime64[us]} passes \texttt{datetime64[ns]}, \texttt{boolean} passes \texttt{bool}, \texttt{Float64} passes \texttt{float64}, and \texttt{object} and \texttt{string[python]} pass \texttt{str}, so the declared dtype stays \texttt{schema.\allowbreak{}py}'s to check (EDN-31 found the datetime unit gap in Pandera too).
\item \textbf{The fill-rate rule checks preprocessing, not the scrape.} This was true when \#25 merged: the interim suite asserted filled the columns the raw contract declared non-null, and a gap in a scrape failed \texttt{download} first. \href{https://github.com/mlops-2627q1-mds-upc/MLOPS_Recommenditos/issues/78}{\#78} moved those fill rates into the raw suite (EDN-77).
\item \textbf{No vacuous pass.} Every column rule holds over no rows, and Great Expectations skips missing values, so an empty frame or a bounded column that has emptied out would pass every rule that reads it. Both suites therefore check the row count against \texttt{validate.\allowbreak{}min\_\allowbreak{}rows} and the fill rate of \texttt{registration\_\allowbreak{}date} and \texttt{mileage\_\allowbreak{}km\_\allowbreak{}raw} against \texttt{validate.\allowbreak{}filled\_\allowbreak{}mostly}, and \texttt{validate-\allowbreak{}data} refuses a suite with no expectations and a run with fewer results than expectations.
\item \textbf{The raw date rule on the parsed date.} Great Expectations types a bound as a number or a date: a text bound is parsed into a \texttt{date}, and comparing it with the raw file's text column raises. So the raw suite sees the published file with \texttt{registration\_\allowbreak{}date} cast by the interim contract's own cast, and every other column untouched. It therefore flags exactly the rows the interim rule then removes.
\item \textbf{What was taken out of \#25's list, and why.} An \texttt{is\_\allowbreak{}used} / \texttt{offer\_\allowbreak{}type} agreement rule would fail 18,108 of the published file's 113,708 used passenger-car rows by design (EDN-23). The supported-make list cannot be a rule here: \texttt{split} computes it after this stage, and the processed sets deliberately keep every make (EDN-48), so no frame \texttt{validate-\allowbreak{}data} sees is meant to satisfy it. Its guarantee is enforced where the list is applied, in \texttt{features} (EDN-67), and checked again by \texttt{train} and \texttt{evaluate}.
\item \textbf{How a failure gates.} Any failed expectation fails the stage, as the structural contract already did, and every frame's result is read; the demo only logs a count, and only for the first validation result of its checkpoint run. The summary and the Data Docs are written before the raise, so a failed run says which rule broke and shows a sample of the failing values.
\end{itemize}
\par
On the real snapshot every expectation of both suites passes. The raw suite's date rule finds 164 registrations after the reference date, 0.1388\,\% of the 118,161 dated rows and so inside the 1\,\% tolerance, and its mileage rule three readings above 1,000,000 km, which EDN-72 is about.
\par
One cost of the tool was found only on the real data: Great Expectations hashes every in-memory batch into a fingerprint for the result's markers, which on the raw frame doubles the stage's memory. Measured with \texttt{reports/\allowbreak{}analysis/\allowbreak{}gx\_\allowbreak{}fingerprint\_\allowbreak{}cost.\allowbreak{}py}, the stage takes 14 s and peaks at 4.4 GB with the fingerprint, and 11 s and 2.2 GB without it, the raw frame itself being most of that. \texttt{validate-\allowbreak{}data} switches the fingerprint off for its own runs, and a test fails if an upgrade makes that switch stop working.},
  ai = {Information seeking, Alternative generation, Alternative assessment, Recommendation, Solution generation},
  response = {Accepted with modifications},
  assessment = {AI implemented \#25 and proposed every choice above, with the alternatives, after testing three assumptions before building on them, each of which changed the design. Running every expectation kind on a satisfying and a breaking frame showed that Great Expectations works on pandas 3 but cannot bound a text date. Measuring what the store writes across two builds turned ``commit \texttt{gx/\allowbreak{}}'', the demo's layout and the stub's own plan, into a cached output. Checking each listed rule against the point in the DAG where it runs found that the supported-make rule could not mean anything there, which the issue, \texttt{schema.\allowbreak{}py} and \texttt{docs/\allowbreak{}docs/\allowbreak{}pipeline.\allowbreak{}md} all still assumed. Lukas confirmed the setup as proposed. An independent AI review of PR \#76 then found four things the first version got wrong, all fixed before merge: the fill-rate rule was described as catching a scrape gap that the raw contract already stops at \texttt{download}; a zero-row frame, an emptied bounded column or an empty suite passed vacuously; the type blind spots were understated as two, while \texttt{object}, \texttt{string[python]} and \texttt{Float64} pass as well; and the generated not-null checks had no test, so deleting them passed the suite.},
  aievidence = {Claude Code session on 2026-10-05 implementing issue \#25: throwaway probes on the fixture and the real snapshot before any stage code, the store measurement, and the options above presented with a recommendation, which Lukas confirmed; then an independent AI review of PR \#76, whose findings are listed above.},
  evidence = {\href{https://github.com/mlops-2627q1-mds-upc/MLOPS_Recommenditos/blob/main/recommenditos/data/gx_context_configuration.py}{\texttt{recommenditos/\allowbreak{}data/\allowbreak{}gx\_\allowbreak{}context\_\allowbreak{}configuration.\allowbreak{}py}} and \href{https://github.com/mlops-2627q1-mds-upc/MLOPS_Recommenditos/blob/main/recommenditos/data/validate_data.py}{\texttt{recommenditos/\allowbreak{}data/\allowbreak{}validate\_\allowbreak{}data.\allowbreak{}py}}, whose docstrings give each rule's reason; \href{https://github.com/mlops-2627q1-mds-upc/MLOPS_Recommenditos/blob/main/tests/test_validate_data.py}{\texttt{tests/\allowbreak{}test\_\allowbreak{}validate\_\allowbreak{}data.\allowbreak{}py}}; \href{https://github.com/mlops-2627q1-mds-upc/MLOPS_Recommenditos/blob/main/reports/analysis/gx_probe.py}{\texttt{reports/\allowbreak{}analysis/\allowbreak{}gx\_\allowbreak{}probe.\allowbreak{}py}} and \href{https://github.com/mlops-2627q1-mds-upc/MLOPS_Recommenditos/blob/main/reports/analysis/gx_fingerprint_cost.py}{\texttt{gx\_\allowbreak{}fingerprint\_\allowbreak{}cost.\allowbreak{}py}} with their results; \href{https://github.com/mlops-2627q1-mds-upc/MLOPS_Recommenditos/blob/main/reports/data-validation/summary.json}{\texttt{reports/\allowbreak{}data-\allowbreak{}validation/\allowbreak{}summary.\allowbreak{}json}}; \href{https://github.com/mlops-2627q1-mds-upc/MLOPS_Recommenditos/blob/main/docs/docs/pipeline.md}{pipeline docs}, ``Data validation''; \hyperref[EDN-22]{EDN-22}, \hyperref[EDN-23]{EDN-23}, \hyperref[EDN-31]{EDN-31}, \hyperref[EDN-48]{EDN-48}; \href{https://github.com/mlops-2627q1-mds-upc/MLOPS_Recommenditos/issues/25}{issue \#25}, \href{https://github.com/mlops-2627q1-mds-upc/MLOPS_Recommenditos/pull/76}{PR \#76}.},
}
